\documentclass[11pt,a4paper]{article}
\usepackage[utf8]{inputenc}
\usepackage[T1]{fontenc}
\usepackage{geometry}
\usepackage{hyperref}
\geometry{margin=0.8in}
\setlength{\parindent}{0pt}
\setlength{\parskip}{6pt}
\begin{document}
\title{cap-com-video v1}
\date{}
\maketitle
Runnable release bundle for video review, CRF quality testing, and image/video sharpness comparison in the \texttt{non-flip} / \texttt{flip} workflow.
Default release path:
\begin{verbatim}
/home/cra-space-center/Desktop/real-6-day-26-30-copy/cap-com-video/
\end{verbatim}
\section{Release Contents}
\begin{verbatim}
cap-com-video/
  README.md
  README.tex
  README.pdf
  requirements.txt
  environment.yml
  video_review_center.py
  test_video_crop_calibrator_crf.py
  compare_image_sharpness.py
  build_clip_edge_frames.py
  build_long_video_from_clip_ranges.py
  crop_calibration_space_24.json
\end{verbatim}
\begin{itemize}
\item \texttt{video\_review\_center.py} starts a local browser UI for reviewing processed \texttt{flip} videos, optional \texttt{non-flip} sources, event/session data, long-video jobs, and edge-frame jobs.
\item \texttt{test\_video\_crop\_calibrator\_crf.py} renders one calibrated input video at multiple CRF values so you can choose the best size/quality balance.
\item \texttt{compare\_image\_sharpness.py} compares image or video inputs and writes side-by-side preview images plus a JSON metrics report.
\item \texttt{build\_clip\_edge\_frames.py} and \texttt{build\_long\_video\_from\_clip\_ranges.py} are runtime helpers imported by \texttt{video\_review\_center.py}.
\item \texttt{crop\_calibration\_space\_24.json} is a sample calibration profile for CRF testing.
\end{itemize}
\section{System Requirements}
Install system packages:
\begin{verbatim}
sudo apt update
sudo apt install -y python3 ffmpeg
\end{verbatim}
\texttt{ffmpeg} and \texttt{ffprobe} must be available in \texttt{PATH}. They are used for rendering, probing duration/size, building long videos, and extracting frames.
\section{Conda Setup}
Create and activate the conda environment:
\begin{verbatim}
cd /home/cra-space-center/Desktop/real-6-day-26-30-copy/cap-com-video
conda env create -f environment.yml
conda activate cap-com-video
\end{verbatim}
If the environment already exists, update it:
\begin{verbatim}
conda env update -n cap-com-video -f environment.yml --prune
conda activate cap-com-video
\end{verbatim}
Alternative pip install inside an existing conda environment:
\begin{verbatim}
conda activate vdobrowser
python -m pip install -r requirements.txt
\end{verbatim}
Check external tools:
\begin{verbatim}
ffmpeg -version
ffprobe -version
\end{verbatim}
\section{1. Start Video Review Center}
Run the local web server:
\begin{verbatim}
cd /home/cra-space-center/Desktop/real-6-day-26-30-copy/cap-com-video
conda activate cap-com-video
python video_review_center.py \
  --flip /home/cra-space-center/Desktop/real-6-day-26-30-copy/flip \
  --non-flip /home/cra-space-center/Desktop/real-6-day-26-30-copy/non-flip \
  --host 127.0.0.1 \
  --port 8765
\end{verbatim}
Open:
\begin{verbatim}
http://127.0.0.1:8765/
\end{verbatim}
Useful pages:
\begin{verbatim}
http://127.0.0.1:8765/                 video review UI
http://127.0.0.1:8765/events           event/session map
http://127.0.0.1:8765/events-export.html
http://127.0.0.1:8765/long-builder.html
http://127.0.0.1:8765/edge-frame-builder.html
\end{verbatim}
Important options:
\begin{verbatim}
python video_review_center.py --help
\end{verbatim}
\begin{itemize}
\item \texttt{--flip}: root folder for processed videos.
\item \texttt{--non-flip}: optional root folder for original videos.
\item \texttt{--command-csv}: command export CSV/TSV for event/session data.
\item \texttt{--remote-video-csv}: remote video list CSV for matching remote clips.
\item \texttt{--event-notes}: JSON file used to save event group names and notes.
\item \texttt{--refresh-seconds}: auto-refresh interval for the video index when API requests arrive. Use \texttt{0} to disable.
\end{itemize}
Stop the server with \texttt{Ctrl+C}.
\section{2. Test CRF Quality}
Use this tool to choose a CRF value that keeps enough image detail without producing unnecessarily large files:
\begin{verbatim}
cd /home/cra-space-center/Desktop/real-6-day-26-30-copy/cap-com-video
conda activate cap-com-video
python test_video_crop_calibrator_crf.py /path/to/input.mp4 \
  --config ./crop_calibration_space_24.json \
  --output-dir ./crf_test \
  --crf-values 12 16 18 20 23 \
  --max-seconds 15 \
  --sample-seconds 5 \
  --overwrite
\end{verbatim}
CRF guide:
\begin{verbatim}
Lower CRF = higher quality, sharper detail, larger files
Higher CRF = stronger compression, smaller files, more softness/artifacts
CRF 12-16 = high quality for detail inspection
CRF 18-20 = good starting range for v1 balancing
CRF 23+   = smaller files, but visual checking is required
\end{verbatim}
Outputs:
\begin{verbatim}
crf_test/
  input_stem_crf12.mp4
  input_stem_crf12_sample.png
  input_stem_crf16.mp4
  input_stem_crf16_sample.png
  ...
  summary.json
  note.txt
\end{verbatim}
\begin{itemize}
\item \texttt{summary.json} stores output path, duration, size, and bitrate for each CRF.
\item \texttt{note.txt} records the input, config, preset, max seconds, sample seconds, and CRF guide.
\item \texttt{*\_sample.png} files are quick side-by-side inspection frames for texture, edges, text, fine detail, and noise.
\end{itemize}
\section{3. Compare Sharpness}
Compare PNG frames from a CRF test:
\begin{verbatim}
cd /home/cra-space-center/Desktop/real-6-day-26-30-copy/cap-com-video
conda activate cap-com-video
python compare_image_sharpness.py \
  --input ./crf_test/input_stem_crf16_sample.png \
  --input ./crf_test/input_stem_crf18_sample.png \
  --label CRF16 \
  --label CRF18 \
  --output-dir ./sharpness_compare
\end{verbatim}
Compare frames directly from videos:
\begin{verbatim}
python compare_image_sharpness.py \
  --input /path/to/raw.mp4 \
  --input /path/to/calibrated.mp4 \
  --label raw \
  --label calibrated \
  --time-seconds 10 \
  --output-dir ./sharpness_compare
\end{verbatim}
Compare a specific region of interest:
\begin{verbatim}
python compare_image_sharpness.py \
  --input ./crf_test/input_stem_crf16_sample.png \
  --input ./crf_test/input_stem_crf20_sample.png \
  --label CRF16 \
  --label CRF20 \
  --crop 100,80,400,240 \
  --output-dir ./sharpness_compare_roi
\end{verbatim}
Outputs:
\begin{verbatim}
sharpness_compare/
  comparison_strip.jpg
  comparison_zoom_strip.jpg
  report.json
\end{verbatim}
\begin{itemize}
\item \texttt{comparison\_strip.jpg} is a side-by-side preview with key metrics.
\item \texttt{comparison\_zoom\_strip.jpg} is a centered zoom crop for detail/noise inspection.
\item \texttt{report.json} includes metrics such as \texttt{laplacian\_var}, \texttt{tenengrad}, \texttt{brenner}, \texttt{mean\_gradient}, \texttt{noise\_std}, \texttt{psnr}, and \texttt{ssim}.
\end{itemize}
Metrics are screening tools, not final visual judgment. Noisy footage can score as sharp even when it does not look better.
\section{v1 Release Checklist}
Before marking this bundle as v1-ready:
\begin{verbatim}
cd /home/cra-space-center/Desktop/real-6-day-26-30-copy/cap-com-video
conda activate cap-com-video
python -m py_compile \
  video_review_center.py \
  test_video_crop_calibrator_crf.py \
  compare_image_sharpness.py \
  build_clip_edge_frames.py \
  build_long_video_from_clip_ranges.py

python video_review_center.py --help
python test_video_crop_calibrator_crf.py --help
python compare_image_sharpness.py --help
\end{verbatim}
Validation with real media:
\begin{verbatim}
1. Video Review Center starts and indexes more than 0 videos.
2. http://127.0.0.1:8765/ opens and can play videos.
3. CRF testing completes on a short sample video and writes summary.json/note.txt.
4. Sharpness comparison completes with at least two PNG or MP4 inputs.
5. comparison_strip.jpg and comparison_zoom_strip.jpg open correctly and labels are readable.
\end{verbatim}
If the review UI opens but event/session data is empty, verify \texttt{--command-csv} and \texttt{--remote-video-csv}.
\end{document}
